% DAFTAR SIMBOL
\chapter*{DAFTAR SIMBOL}
\addcontentsline{toc}{chapter}{Daftar Simbol}

Berikut adalah daftar simbol dan notasi yang digunakan dalam penelitian ini:

\vspace{0.5cm}

\noindent
\begin{tabularx}{\textwidth}{@{} l X @{}}
\toprule
\textbf{Simbol} & \textbf{Keterangan} \\
\midrule

\multicolumn{2}{l}{\textbf{A. Simbol Matematika Dasar}} \\
$\alpha, \beta, \gamma$ & Huruf Yunani (alpha, beta, gamma) \\
$\pi$ & Konstanta pi ($\approx 3.14159$) \\
$e$ & Konstanta euler ($\approx 2.71828$) \\
$\infty$ & Tak terhingga (infinity) \\
$\emptyset$ & Himpunan kosong (empty set) \\
$\in, \notin$ & Elemen dari, bukan elemen dari (element, not element) \\
$\subset, \subseteq$ & Himpunan bagian, bagian atau sama dengan \\
$\cup, \cap$ & Gabungan, irisan (union, intersection) \\
$\forall$ & Untuk semua (for all) \\
$\exists$ & Ada/terdapat (exists) \\
$\Rightarrow, \Leftrightarrow$ & Implikasi, biimplikasi (implication, biconditional) \\

\multicolumn{2}{l}{\textbf{B. Simbol Operasi}} \\
$+, -, \times, \div$ & Penjumlahan, pengurangan, perkalian, pembagian \\
$a^b$ & $a$ pangkat $b$ (exponent) \\
$\sqrt{a}$ & Akar kuadrat dari $a$ (square root) \\
$|a|$ & Nilai mutlak dari $a$ (absolute value) \\
$\sum_{i=1}^{n}$ & Penjumlahan (summation) dari $i$ ke $n$ \\
$\prod_{i=1}^{n}$ & Perkalian (product) dari $i$ ke $n$ \\
$\int_a^b$ & Integral dari $a$ ke $b$ (definite integral) \\
$\frac{\partial}{\partial x}$ & Turunan parsial terhadap $x$ (partial derivative) \\

\multicolumn{2}{l}{\textbf{C. Simbol Perbandingan}} \\
$=$, $\neq$ & Sama dengan, tidak sama dengan (equal, not equal) \\
$<, >$ & Kurang dari, lebih dari (less than, greater than) \\
$\leq, \geq$ & Kurang dari atau sama dengan, lebih dari atau sama dengan \\
$\approx$ & Kira-kira sama dengan (approximately equal) \\
$\propto$ & Berbanding lurus dengan (proportional to) \\

\multicolumn{2}{l}{\textbf{D. Simbol Aljabar Linear}} \\
$\mathbf{A}, \mathbf{B}$ & Matriks (matrices) \\
$\mathbf{a}, \mathbf{b}$ & Vektor (vectors) \\
$A^T$ & Transpose matriks $A$ \\
$A^{-1}$ & Invers matriks $A$ \\
$\det(A)$ & Determinan matriks $A$ \\
$\text{tr}(A)$ & Trace (jejak) matriks $A$ \\
$\mathbf{a} \cdot \mathbf{b}$ & Dot product (perkalian skalar) \\
$\mathbf{a} \times \mathbf{b}$ & Cross product (perkalian silang) \\
$\|\mathbf{a}\|$ & Norma/panjang vektor $\mathbf{a}$ \\

\multicolumn{2}{l}{\textbf{E. Simbol Fungsi}} \\
$f(x)$ & Fungsi $f$ dengan variabel $x$ \\
$f'(x), f''(x)$ & Turunan pertama dan kedua dari $f(x)$ \\
$\frac{df}{dx}$ & Turunan $f$ terhadap $x$ \\
$\lim_{x \to a}$ & Limit saat $x$ mendekati $a$ \\
$\sin, \cos, \tan$ & Sinus, cosinus, tangen (trigonometric) \\
$\log, \ln$ & Logaritma dasar 10, logaritma natural \\
$\exp(x)$ & Exponential function, $e^x$ \\

\multicolumn{2}{l}{\textbf{F. Simbol Probabilitas dan Statistika}} \\
$P(A)$ & Probabilitas kejadian $A$ \\
$E[X]$ & Nilai ekspektasi/harapan dari $X$ \\
$\text{Var}(X)$ & Varians dari $X$ \\
$\sigma$ & Standar deviasi (standard deviation) \\
$\mu$ & Rata-rata/mean \\
$\rho$ & Korelasi (correlation) \\
$\chi^2$ & Chi-squared distribution \\
$\mathcal{N}(\mu, \sigma^2)$ & Distribusi normal dengan mean $\mu$ dan varians $\sigma^2$ \\

\multicolumn{2}{l}{\textbf{G. Simbol Fisika}} \\
$v$ & Kecepatan (velocity) \\
$a$ & Percepatan (acceleration) \\
$F$ & Gaya (force) \\
$E$ & Energi (energy) \\
$P$ & Daya (power) \\
$\lambda$ & Panjang gelombang (wavelength) \\
$\nu$ & Frekuensi (frequency) \\
$h$ & Konstanta Planck \\
$c$ & Kecepatan cahaya \\
$G$ & Konstanta gravitasi \\

\multicolumn{2}{l}{\textbf{H. Simbol Kimia}} \\
$n$ & Jumlah mol \\
$M$ & Massa molar \\
$\Delta H$ & Perubahan entalpi \\
$\Delta G$ & Perubahan energi Gibbs \\
$K_c$ & Konstanta kesetimbangan \\
$pH$ & Derajat keasaman \\
$\lambda$ & Panjang gelombang \\

\multicolumn{2}{l}{\textbf{I. Simbol Umum Lainnya}} \\
$\approx$ & Kira-kira sama dengan \\
$\because$ & Karena (because) \\
$\therefore$ & Oleh karena itu (therefore) \\
$\Box, \square$ & Akhir pembuktian (end of proof) \\
$\dagger, \ddagger$ & Catatan kaki/superskrip \\
$\ast$ & Operasi konvolusi (convolution) \\
$\nabla$ & Operator nabla/gradient \\
$\partial$ & Turunan parsial (partial derivative) \\

\bottomrule
\end{tabularx}

\vspace{0.5cm}

\noindent
\textbf{Catatan:}
\begin{itemize}
  \item Simbol-simbol di atas adalah simbol standar LaTeX yang umum digunakan dalam penelitian akademik.
  \item Jika ada simbol khusus dalam penelitian Anda, tambahkan ke daftar ini dengan penjelasan yang jelas.
  \item Pastikan penggunaan simbol konsisten di seluruh dokumen.
  \item Untuk simbol yang tidak ada di daftar ini, gunakan paket \texttt{amssymb} atau \texttt{amsmath} di LaTeX.
\end{itemize}
